% !TeX root = ../T_tuner_control_analysis.tex
% SNL tuner loop state machine.  loop_off and loop_reset sit side by side
% because loop_unknown branches to one or the other; both rejoin at loop_on.
\begin{tikzpicture}[font=\scriptsize, x=1mm, y=1mm]
\def\xC{56}     % centre column

\node[lbl,anchor=south] at (\xC,6) {startup};
\node[blk,text width=54mm,anchor=north] (init) at (\xC,0)
  {\textbf{\texttt{loop\_init}}\\{\tiny initialise PV monitors, wait for valid readbacks}};
\draw[sig] (\xC,5) -- (init.north);

\node[blk,text width=54mm,anchor=north] (unk) at (\xC,-22)
  {\textbf{\texttt{loop\_unknown}}\\{\tiny wait for the loop control to be set}};
\draw[sig] (init.south) -- node[right,tag,xshift=1mm]
  {\texttt{sm\_dmov} and \texttt{sm\_posn} valid} (unk.north);

% the branch: idle or reset
\def\yB{-48}
\draw[draw=cNeut,thick] (unk.south) -- (\xC,\yB);
\draw[draw=cNeut,thick] (22,\yB) -- (90,\yB);
\node[blk,text width=40mm,anchor=north] (off) at (22,\yB-4)
  {\textbf{\texttt{loop\_off}}};
\node[blk,text width=40mm,anchor=north] (rst) at (90,\yB-4)
  {\textbf{\texttt{loop\_reset}}\\{\tiny return the motor to home}};
\draw[sig] (22,\yB) -- (off.north);
\draw[sig] (90,\yB) -- (rst.north);
\node[tag,anchor=south east] at (20,\yB+1) {ctrl $=$ IDLE};
\node[tag,anchor=south west] at (92,\yB+1) {ctrl $=$ RESET};

% both rejoin at loop_on; start the drops at the box edges, not a guessed y
\def\yJ{-74}
\draw[draw=cNeut,thick] (off.south) -- (22,\yJ) -- (90,\yJ);
\draw[draw=cNeut,thick] (rst.south) -- (90,\yJ);
\node[tag,anchor=west] at (92,-70) {reset complete};

\node[ctrl,text width=76mm,anchor=north] (on) at (\xC,\yJ-4)
  {\textbf{\texttt{loop\_on}}\\{\tiny main control state: measure, compute delta,
   command the move, verify the position}};
\draw[sig] (\xC,\yJ) -- (on.north);
\end{tikzpicture}
